% Candidate research section -- mathematical manuscript proof, not yet Lean-certified.
% It may be incorporated after independent peer review and a CI-verified formalization.
\subsection{Finite-horizon transport through conditioned shell histories}
\label{subsec:finite-horizon-transport}
Let $c\ge-1$ be odd and let
\[
 T_c(n)=\begin{cases}n/2,&n\equiv0\pmod2,\\(3n+c)/2,&n\equiv1\pmod2.\end{cases}
\]
Fix positive integers $a,L$, a starting interval $J=[A,B)\cap\mathbb N$
with $1\le A<B$, and the shell $I_L=(L/2,L]\cap\mathbb N$.
An orbit is stopped immediately when its state is at most $a$.
For integers $k\ge1$ and $t\ge0$, let $E_{k,t}\subset J$ consist
of the starts for which the $k$-th visit to $I_L$ (counting the
initial state, if applicable) occurs exactly at time $t$ before stopping.
Set $E_{k,\le T}=\bigsqcup_{t=0}^T E_{k,t}$ and
\[
 G_c=\{y\in\mathbb N:y\equiv0\pmod4\text{ or }
 3y+c\equiv0\pmod{16}\}.
\]
Because $c$ is odd, $G_c$ consists of the five distinct classes
$0,4,8,12,5c\pmod{16}$.

\begin{proposition}[Finite-horizon conditional modulo-16 transport]
\label{prop:finite-horizon-transport}
For every $c,a,L,A,B,k,T$ as above, uniformly in the values of
these parameters,
\begin{align}
 \left|\sum_{n\in E_{k,\le T}}
  \left(\mathbf1_{G_c}(T_c^{\tau_k(n)}n)-\frac5{16}\right)\right|
  &\le40(T+1)2^T,\label{eq:finite-horizon-count}\\
 \left|\sum_{n\in E_{k,\le T}}\frac1n
  \left(\mathbf1_{G_c}(T_c^{\tau_k(n)}n)-\frac5{16}\right)\right|
  &\le\frac{40(T+1)2^T}{A}.\label{eq:finite-horizon-log}
\end{align}
Here $\tau_k(n)$ is the time of the $k$-th shell visit, so that
$\tau_k(n)\le T$ on $E_{k,\le T}$. In particular, if the total
weighted mass $W_{k,T}=\sum_{n\in E_{k,\le T}}1/n$ is positive, then
\[
 \left|\frac{\sum_{n\in E_{k,\le T}}
                n^{-1}\mathbf1_{G_c}(T_c^{\tau_k(n)}n)}{W_{k,T}}
      -\frac5{16}\right|
 \le\frac{40(T+1)2^T}{A W_{k,T}}.
\]
\end{proposition}
\begin{proof}
Fix $t\le T$. For each residue $r\pmod{2^t}$, the first $t$
parities of $T_c$ are fixed on the progression $n=r+2^t m$.
If $h_s$ is the number of odd steps among the first $s$ steps,
induction gives, for $0\le s\le t$,
\[
 T_c^s(r+2^t m)=T_c^s(r)+3^{h_s}2^{t-s}m.
\]
All these functions are strictly increasing affine functions of $m$.
In particular the state at time $t$ advances by the \emph{odd} step
$3^{h_t}$ as $m$ increases by one. Thus its residues modulo $16$
traverse all $16$ classes in each consecutive block of $16$ values
of $m$, with exactly five values in $G_c$.

To determine whether the $k$-th shell visit occurs at time $t$
before the cutoff, it is enough to know, for each $0\le s\le t$,
on which side of the three thresholds $a,L/2,L$ the state
$T_c^s(n)$ lies. Because each state is increasing affine in $m$,
these comparisons partition the admissible $m$-interval (from $J$)
into at most $3(t+1)+1$ consecutive integer subintervals. On each
subinterval the full stopping/shell-visit profile, and hence
membership of $n$ in $E_{k,t}$, is constant. Over any one such
subinterval the discrepancy from the expected $5/16$ count is
at most $5$: complete blocks of $16$ cancel, leaving at most
$15$ terms. Hence the discrepancy for a fixed residue class
$r\pmod{2^t}$ is at most $5[3(t+1)+1]$. There are $2^t$
residue classes, so the discrepancy for $E_{k,t}$ is at most
$5(3t+4)2^t$.

For the harmonic weights $1/n$, the same bound holds after
multiplication by $1/A$. Indeed, on each consecutive subinterval
of $m$ the partial sums of the centred $16$-periodic sequence
$\mathbf1_{G_c}(T_c^t(r+2^t m))-5/16$ are bounded in absolute
value by $5$. Summation by parts against the positive,
nonincreasing weights $(r+2^t m)^{-1}$ bounds the weighted
contribution of that subinterval by $5/A$.

Finally, the events $E_{k,t}$ for different $t$ are disjoint.
Summing the two bounds over $0\le t\le T$ gives at most
\[
 5(3T+4)\sum_{t=0}^T2^t
 \le 10(3T+4)2^T\le40(T+1)2^T,
\]
which proves both assertions.
\end{proof}

\begin{corollary}[A controlled fraction of logarithmic time]
\label{cor:finite-horizon-transport}
Suppose $A\to\infty$, the starting intervals have $B/A$
bounded above and away from $1$, and $T=T(A)$ satisfies
$T\le(1-\eta)\log_2 A$ for some fixed $\eta>0$. Then the
absolute errors in \eqref{eq:finite-horizon-log} vanish as
$O((\log A)A^{-\eta})$. In particular, on any sequence
of events $E_{k,\le T}$ of harmonic mass at least
$\delta>0$, the \emph{conditional} probability of $G_c$
converges to $5/16$. The same holds for uniform counting
when $\#E_{k,\le T}$ is bounded below by a fixed positive
multiple of $A$.
\end{corollary}


\begin{corollary}[Early--late reduction for the full entrance law]
\label{cor:finite-horizon-early-late}
Let $E_{k,<\infty}=\bigcup_{t\ge0}E_{k,t}$ (a disjoint union),
$W=\sum_{n\in E_{k,<\infty}}1/n>0$, and
$W_{>T}=\sum_{n\in E_{k,<\infty}\setminus E_{k,\le T}}1/n$.
Then
\[
 \mathbb P_{\log}(Y_k\in G_c\mid E_{k,<\infty})
 \ge \frac5{16}\left(1-\frac{W_{>T}}W\right)
     -\frac{40(T+1)2^T}{AW},
\]
where $\mathbb P_{\log}$ is the finite harmonic probability on $J$,
conditioned on $E_{k,<\infty}$.
\end{corollary}
\begin{proof}
Apply the lower half of \eqref{eq:finite-horizon-log} to the early
entrances and use the nonnegativity of the good-class mass on late
entrances. Divide by $W$.
\end{proof}



\begin{proposition}[Unrestricted-time absolute transport at a shallow cutoff]
\label{prop:full-time-shallow-mod16}
Fix an odd $c\ge-1$, let $X\to\infty$ through integers, and
choose any cutoff $a_X$ satisfying
$X^{19/20}\le a_X<X/2$. Equip either all integers or just the
odd integers in $J_X=[4X,16X)\cap\mathbb N$ with their normalized
harmonic measure $\mu_X$. For the $k$-th visit $Y_k$ to
$I_X=(X/2,X]$ \emph{before} the first passage at or below $a_X$,
write $E_k$ for the event that this visit occurs (allowing arbitrarily
many iteration steps). There is an absolute constant
$\gamma=1/(1024\log 2)>0$ such that, for every fixed $c$,
\[
 \sup_{k\ge1}\left|
 \mu_X(E_k\cap\{Y_k\in G_c\})-\frac5{16}\mu_X(E_k)
 \right|=O_c(X^{-\gamma}).
\]
Consequently, whenever $\mu_X(E_k)\ge\delta>0$, the conditional
probability of $G_c$ equals $5/16+O_{c,\delta}(X^{-\gamma})$.
This estimate is uniform over $k$ in its \emph{absolute} form;
it is not a relative estimate when $\mu_X(E_k)$ is small.
\end{proposition}
\begin{proof}
Put $t=\lfloor\tfrac12\log_2 X\rfloor$. A parity word of
length $t$ has $K$ odd steps and gives
\[
 T_c^t(n)=\frac{3^K}{2^t}n+b_t,\qquad
 b_t\le c_+\bigl((3/2)^t-1\bigr),\quad c_+=\max(c,0).
\]
The second inequality follows by iterating the affine maps
$n\mapsto n/2$ and $n\mapsto(3n+c)/2$ and bounding each
positive additive contribution by a geometric series.
If $K\le9t/16$ and $n<16X$, then
\[
 T_c^t(n)\le16X(3^{9/16}/2)^t
     +c_+((3/2)^t-1)
    =O_c\bigl(X^{1-\delta_0}+X^{\frac12\log_2(3/2)}\bigr),
\]
where
$\delta_0=\frac12(1-\frac9{16}\log_2 3)>1/20$.
Thus, for sufficiently large $X$, every such orbit has passed
below $a_X\ge X^{19/20}$ by time $t$.

Parity words of length $t$ are in bijection with residues modulo
$2^t$. Over $J_X$ their frequencies differ from uniform by an
additive $O(2^t/X)$; the same holds when restricting to odd
starts, with the first bit then fixed to $1$. Under the uniform
word law, $K$ is $\operatorname{Bin}(t,1/2)$ for unrestricted
starts and $1+\operatorname{Bin}(t-1,1/2)$ for odd starts.
Hoeffding's inequality gives
\[
 \Pr(K>9t/16)\le e^{-t/128}
 \quad\text{or}\quad e^{-t/512},
\]
respectively, for sufficiently large $t$. Since all harmonic
weights in $J_X$ differ by a factor at most $4$, this yields
\[
 \mu_X\{\text{first passage below }a_X\text{ takes more
 than }t\text{ steps}\}
 =O_c(X^{-\gamma}).
\]
Apply Proposition~\ref{prop:finite-horizon-transport} to the
$k$-th shell visit up to time $t$. Its count error is
$O(t2^t)$ and its harmonic error is $O(t2^t/X)$.
The same bound holds for odd-only starts (the relevant times
$t'\ge1$ have parity cylinders with odd first bit).
The harmonic mass of $J_X$ is bounded below by a positive
constant for either choice of starting parity, so the normalized
early-event error is $O((\log X)X^{-1/2})$.
The unrestricted event $E_k$ differs from this early event
only on starts not yet past the cutoff by time $t$, whose
harmonic probability is $O_c(X^{-\gamma})$. Because the
centered indicator is bounded by $1$, the last two estimates
prove the claim uniformly in $k$.
\end{proof}


\begin{corollary}[A polynomial starting window matching the manuscript scales]
\label{cor:polynomial-starts-mod16}
Put $\beta=1/200$, and for sufficiently large $X$ take odd starting
values in
\[
 J_X=\bigl[\lceil X^{1+\beta}\rceil,
             \lfloor X^{(1+\beta)^2}\rfloor\bigr],
 \qquad a_X=\lfloor X^{1-2\beta}\rfloor.
\]
For the accelerated map $T_c$, the uniform-in-$k$ \emph{absolute}
estimate of Proposition~\ref{prop:full-time-shallow-mod16} remains valid,
with the same exponent $\gamma$ and the shell $I_X=(X/2,X]$.
This statement concerns accelerated-map shell visits; it does not
identify them with the block observables in the Syracuse-profile theorem.
\end{corollary}
\begin{proof}
Use $t=\lfloor\frac12\log_2X\rfloor$ again. On a parity word
with $K\le9t/16$ the preceding estimate gives
\[
 T_c^t(n)\ll_c X^{(1+\beta)^2-\delta_0}
              +X^{\frac12\log_2(3/2)}
              =o(X^{1-2\beta}),
\]
since $(1+\beta)^2-\delta_0<1-2\beta$.
The proportion of parity words failing the binomial bound is
$O(X^{-\gamma})$. For the odd harmonic starting law on $J_X$,
the same bound holds after an $O(2^t/\min J_X)$ error in
unnormalized harmonic mass: the bad-word indicator is periodic
modulo $2^t$, and summation by parts controls its harmonic
mean. Since the total harmonic mass of odd starts is asymptotic
to $\tfrac12[(1+\beta)^2-(1+\beta)]\log X$, normalization
only improves this error. Apply the finite-horizon harmonic
bound~\eqref{eq:finite-horizon-log}, again observing that
shell visits cannot occur at time zero.
\end{proof}

\paragraph{What this does not prove.}
Proposition~\ref{prop:finite-horizon-transport} applies to \emph{finite-horizon}
shell histories. Proposition~\ref{prop:full-time-shallow-mod16}
extends the \emph{absolute} discrepancy to unrestricted hitting times
for starting windows proportional to $X$ and a shallow polynomial
cutoff. Neither proposition gives relative control
when the conditioning event has very small mass, The polynomial starting windows of the paper are handled by
Corollary~\ref{cor:polynomial-starts-mod16}, but the
\emph{Syracuse block observable} still requires a separate
identification or transport argument.
Corollary~\ref{cor:finite-horizon-early-late} isolates the missing
\emph{late-entrance mass}. In particular, the absolute discrepancy
bound alone cannot furnish a return-index-uniform lower bound on
conditional good-class probability, because the conditioning mass
may be arbitrarily small for large $k$.
It also does not bound the probability of rebounds after the
local descent to $L/4$. Consequently it does not establish the
scale-uniform geometric return condition or $T_X=o(X)$.
The arithmetic of parity cylinders is classical; the point of
the proposition is the explicit uniform discrepancy bound under
shell-history conditioning.
